\scanpage{159}
Pour tout $\unsure{V|W}\in C'^{k}$ ($k$-module de type fini)
soient $K(W)$ \deleted{\unsure{resp.|et}} $K'(W)$ l'ens des ``structures
de type $K$'' resp$^{\mathrm t}$ des structures de $\underline U$-modules
permis sur $k$. On a défini un hom
fonctoriel $K(V)\longrightarrow K'(W)$. \unsure{Ceci|Ceci étant} dit

$1^\circ)$ On veut que cet hom. soit bijectif

$2^\circ)$ Si $V,W\in K$, on veut que les
$K$-hom de $V$ dans $W$ soient identiques aux
$\underline U$-homomorphismes.

On est ramené à la considération des $\underline U$-\deleted{\unsure{lin}}\inserted{modules}
permis qui sont à engendrement fini \deleted{sur} $k$.

$\underline U$ est un \uline{anneau topologique} \inserted{\uline{linéaire, complet avec $1$}\deleted{\unsure{,0}}} ayant un
système \deleted{\unsure{compl}}\inserted{fond.} de \deleted{\unsure{représentations}} \inserted{voisinages $\mathcal J$ \unsure{sts|tels} que}
\deleted{\unsure{de $k$} de $U$ dans $L(V)$, $V$ de type fini sur $k$.} \inserted{$\underline U/\mathcal J$ soit de dim finie.} Tout
anneau de ce type peut s'obtenir ainsi.
\margin{Si $k$ est un corps les $\underline U$ sont les \unsure{duals|duels|duaux} des co-algèbres \unsure{avec} \deleted{\unsure{t}} qui sont réunions \deleted{\unsure{filtrantes de ses}} \deleted{\unsure{algèbres de dim}} de ses co-algèbres de dim finie. Aux hom. continus des $\underline V\to\underline U$ compatibles avec l'augmentation correspondent les hom de Co-Algèbres compatibles avec l'augmentation.}
% slanted

\uline{Première structure supplémentaire}\,: On suppose
\deleted{\unsure{que $k$}} \inserted{tout $V$ \uline{libre} \deleted{\unsure{de $k$}}} puisse être muni canoniquement
d'une structure $\unsure{\mathcal K|K|k}$, i.e. que on ait
sur $k$ une structure privilégiée de $\underline U$ module,
\deleted{\unsure{évidemt}} il revient au $=$ de se donner sur
$\underline U$ une \uline{augmentation continue $\varepsilon$}

% boxed note
\deleted{Les} $\underline U$-modules permis sont les \uline{comodules unitaires}
sur $A$ \deleted{\unsure{qui}} qui sont de dim finie sur $k$.

% archive sheet number: 158
